\section{Concorrência: threads, virtual threads e processos}
\label{sec:vms-concorrencia}

\paragraph{HotSpot.}
A Parte~II mostrou os dois tipos de thread do Java. A thread de plataforma
é uma thread do sistema operacional. A virtual thread
(seção~\ref{sec:virtual-threads}) é escalonada pelo próprio JDK: o
\citetitle{doc:jep444} diz que o escalonador padrão é um
\texttt{ForkJoinPool} em modo FIFO, com paralelismo igual ao número de
processadores disponíveis. Em qualquer dos dois casos, todas as threads
enxergam o mesmo heap: a especificação da JVM diz que ele é
\enquote{shared among all Java Virtual Machine threads}
(\citetitle{doc:jvms25}, §2.5.3). Por isso a Parte~II precisou falar de
\texttt{synchronized} e do modelo de memória
(seções~\ref{sec:sincronizacao} e~\ref{sec:modelo-memoria}).

\paragraph{BEAM.}
A unidade de concorrência do Erlang é o processo, que não é um processo
do sistema operacional. A documentação diz que eles são
\enquote{lightweight (grow and shrink dynamically) with a small memory
footprint} e se comunicam por sinais assíncronos; uma mensagem recebida
vai para o fim da fila do processo, a sua
\termo{mailbox}{Caixa de correio: a fila de mensagens de um processo
Erlang, lida com \texttt{receive}.} (\citetitle{doc:erlang-processes}).
Três números da documentação dão a escala:
\begin{itemize}
\item Um processo recém-criado ocupa 327 palavras de memória, das quais
  233 são o heap inicial (que inclui a pilha). O guia avisa que o número
  veio de uma execução no OTP~27 (erts-14.2.3) e pode variar entre
  versões (\citetitle{doc:erlang-eff-processes}). Em máquina de 64 bits,
  uma palavra tem 8~bytes.
\item O limite padrão de processos simultâneos de um nó é 1\,048\,576
  (opção \texttt{+P} do \citetitle{doc:erlang-erl}).
\item O número de \termo{schedulers}{Escalonadores: threads do sistema
  operacional que executam processos Erlang, cada uma com a sua fila.}
  é, por padrão, o número de processadores lógicos (opção \texttt{+S},
  mesma página).
\end{itemize}
O mesmo guia diz que tudo o que vai numa mensagem entre processos é
\emph{copiado}, com duas exceções no mesmo nó: \emph{refc binaries}
(binários grandes, contados por referência) e literais. Não há memória
compartilhada para proteger com lock: o preço é a cópia.

\paragraph{Preempção por reductions.}
Aqui está a diferença mais funda com as virtual threads. O
\citetitle{livro:stenman} descreve o escalonamento da BEAM como
\enquote{preemptive scheduling on top of cooperative scheduling}: o
processo só pode ser suspenso em certos pontos, como uma chamada de
função ou um \texttt{receive}, mas a VM conta
\termo{reductions}{Unidades de trabalho de um processo Erlang;
grosso modo, cada chamada de função gasta uma. O nome vem do Prolog, em
que cada passo de execução é uma redução.} e, quando o processo gasta a
sua cota, ele é suspenso e outro entra. O livro dá a cota: 4000
reductions (\texttt{CONTEXT\_REDS} em \texttt{erl\_vm.h}); antes do
OTP~20, eram 2000. Como em Erlang não há laço sem chamada de função, nenhum
processo segura um scheduler para sempre.

O JEP~444 diz o contrário para o Java: \enquote{The scheduler does not
currently implement time sharing for virtual threads}. Uma virtual
thread que faz conta sem bloquear fica com a carrier thread até acabar,
o mesmo alerta que o guia do Quarkus repete
(seção~\ref{sec:vt-vs-reativo}).

\paragraph{Onde a BEAM também trava.}
A preempção vale para código Erlang. Uma NIF (função nativa em C chamada
pelo Erlang) roda fora dessa contagem; o \emph{BEAM Book} recomenda
mantê-las curtas \enquote{to avoid blocking the BEAM scheduler} e usar
\emph{dirty schedulers} para operações longas. É o equivalente, na BEAM,
do \emph{pinning} das virtual threads em código nativo
(seção~\ref{sec:virtual-threads}).

\paragraph{O próprio JEP compara.}
O JEP~444 apresenta virtual threads como \emph{user-mode threads}, um
modelo que deu certo em outras linguagens, e cita \enquote{goroutines in
Go and processes in Erlang}. A comparação é justa no custo de criar e
no número de tarefas. Não é justa em dois pontos: o heap continua
compartilhado e não há preempção por tempo. Uma virtual thread é uma
thread barata, não um processo isolado.

\paragraph{GraalVM.}
No modo JIT, o modelo é o da HotSpot, em que a GraalVM se baseia
(\citetitle{doc:graalvm-java}). No Native Image, a linguagem e as APIs de
concorrência são as do Java; o que muda é a compilação e o GC, assunto
das próximas seções.
